\section{Success Criteria}
\label{sec:intro_success_criteria}

To provide a concrete evaluation framework, the following measurable success criteria were defined prior to experimentation, drawing on established conventions in spiking-detection methodology. Each criterion is revisited explicitly in Chapter~\ref{cha:summary_conclusions}.

\begin{enumerate}
    \item \textbf{Detection accuracy parity:}
    On both \gls{miou} and \gls{map_50}, the \gls{snn} must score at least $90\%$ of the \gls{cnn} baseline on the \gls{etram} validation set.
    This $90\%$ threshold is established as a rigorous but realistic target, as achieving exact parity remains an open challenge in deep spiking vision \cite{iaboni2024review}.
    Achieving this will validate that our architecture succeeds in mitigating the compromises expected from the 1-bit temporal quantisation error, across both the geometric-precision metric (\gls{miou}) and the stricter COCO-convention detection-quality metric (\gls{map_50}).

    \item \textbf{Measurable theoretical energy savings:}
    Following the \gls{sop}-based energy methodology derived by \textcite{lemaire2022analytical},
    the \gls{snn} must demonstrate a reduction in theoretical energy consumption per inference relative to the \gls{cnn} baseline, computed from layer-wise \gls{sop} counts and published neuromorphic hardware constants \cite{horowitz20141}.

    \item \textbf{Real-time viability:} 
    Both models must sustain an inference throughput of at least $30$ frames per second on \gls{gpu} hardware, meeting the established minimum threshold for real-time automotive perception \cite{khan2024real}.

    \item \textbf{Architectural parity:} 
    The \gls{snn} and \gls{cnn} must share an equivalent parameter count (within $1\%$), ensuring that any observed performance gap reflects the spiking mechanism rather than differences in model capacity.

    \item \textbf{Detection completeness:}
    Both models must achieve a recall of $100\%$ for objects at the COCO-convention \gls{iou} threshold of $0.5$ \cite{lin2014microsoft} on the validation set, confirming that the multi-box regression head reliably detects all ground-truth vehicles in the scene.

    \item \textbf{Confidence calibration:}
    The \gls{snn} must exhibit confidence calibration equivalent to, or superior to, that of the \gls{cnn} baseline on the \gls{etram} validation set, evaluated via the reliability-diagram framework of~\textcite{guo2017calibration}.
    This criterion is safety-critical: in \gls{adas} applications, a model that systematically over-states its detection confidence may suppress downstream safety interventions.
    A neuromorphic architecture that delivers energy savings at the cost of increased over-confidence would therefore be unsuitable for deployment, regardless of its accuracy or efficiency gains.

\end{enumerate}